\chapter{Of the Embarrassment of Riches}

\lettrine{D}{'Artagnan} lost no time, and as soon as the thing  was suitable and opportune, he paid a visit to the lord treasurer of his  majesty. He had then the satisfaction to exchange a piece of paper, covered  with very ugly writing, for a prodigious number of crowns, recently stamped  with the effigies of his very gracious majesty Charles II.

D'Artagnan easily controlled himself: and yet, on this occasion, he could  not help evincing a joy which the reader will perhaps comprehend, if he deigns  to have some indulgence for a man who, since his birth, had never seen so many  pieces and rolls of pieces juxta-placed in an order truly agreeable to the eye.  The treasurer placed all the rolls in bags, and closed each bag with a stamp  sealed with the arms of England, a favour which treasurers do not grant to  everybody. Then, impassible, and just as polite as he ought to be towards a man  honoured with the friendship of the king, he said to D'Artagnan:

<Take away your money, sir.> Your money! These words made a  thousand chords vibrate in the heart of D'Artagnan, which he had never  felt before. He had the bags packed in a small cart, and returned home  meditating deeply. A man who possessed three hundred thousand livres can no  longer expect to wear a smooth brow; a wrinkle for every hundred thousand  livres is not too much.

D'Artagnan shut himself up, ate no dinner, closed his door to everybody,  and, with a lighted lamp, and a loaded pistol on the table, he watched all  night, ruminating upon the means of preventing these lovely crowns, which from  the coffers of the king had passed into his coffers, from passing from his  coffers into the pockets of any thief whatever. The best means discovered by  the Gascon was to inclose his treasure, for the present, under locks so solid  that no wrist could break them, and so complicated that no master-key could  open them. D'Artagnan remembered that the English are masters in  mechanics and conservative industry; and he determined to go in the morning in  search of a mechanic who would sell him a strong box. He did not go far; Master  Will Jobson, dwelling in Piccadilly, listened to his propositions, comprehended  his wishes, and promised to make him a safety lock that should relieve him from  all future fear.

<I will give you,> said he, <a piece of mechanism entirely  new. At the first serious attempt upon your lock, an invisible plate will open  of itself and vomit forth a pretty copper bullet the weight of a  mark—which will knock down the intruder, and not with a loud report. What  do you think of it?>

<I think it very ingenuous,> cried D'Artagnan; <the  little copper bullet pleases me mightily. So now, sir mechanic, the  terms?>

<A fortnight for the execution, and fifteen hundred livres payable on  delivery,> replied the artisan.

D'Artagnan's brow darkened. A fortnight was delay enough to allow  the thieves of London time to remove all occasion for the strong box. As to the  fifteen hundred livres—that would be paying too dear for what a little  vigilance would procure him for nothing.

<I will think of it,> said he; <thank you, sir.> And he  returned home at full speed; nobody had yet touched his treasure. That same day  Athos paid a visit to his friend and found him so thoughtful that he could not  help expressing his surprise.

<How is this?> said he, <you are rich and not gay—you,  who were so anxious for wealth!>

<My friend, the pleasures to which we are not accustomed oppress us more  than the griefs with which we are familiar. Give me your opinion, if you  please. I can ask you, who have always had money: when we have money, what do  we do with it?>

<That depends.>

<What have you done with yours, seeing that it has not made you a miser  or a prodigal? For avarice dries up the heart, and prodigality drowns  it—is that not so?>

<Fabricius could not have spoken more justly. But in truth, my money has  never been a burden to me.>

<How so? Do you place it out at interest?>

<No; you know I have a tolerably handsome house; and that house composes  the better part of my property.>

<I know it does.>

<So that you can be as rich as I am, and, indeed, more rich, whenever you  like, by the same means.>

<But your rents,—do you lay them by?>

<No.>

<What do you think of a chest concealed in a wall?>

<I never made use of such a thing.>

<Then you must have some confidant, some safe man of business who pays  you interest at a fair rate.>

<Not at all.>

<Good heavens! what do you do with it, then?>

<I spend all I have, and I only have what I spend, my dear  D'Artagnan.>

<Ah! that may be. But you are something of a prince; fifteen or sixteen  thousand livres melt away between your fingers; and then you have expenses and  appearances\longdash>

<Well, I don't see why you should be less of a noble than I am, my  friend; your money would be quite sufficient.>

<Three hundred thousand livres! Two-thirds too much!>

<I beg your pardon—did you not tell me?—I thought I heard you  say—I fancied you had a partner\longdash>

<Ah! Mordioux! that's true,> cried D'Artagnan,  colouring; <there is Planchet. I had forgotten Planchet, upon my life!  Well! there are my three hundred thousand livres broken into. That's a  pity! it was a round sum, and sounded well. That is true, Athos; I am no longer  rich. What a memory you have!>

<Tolerably good; yes, thank God!>

<The worthy Planchet!> grumbled D'Artagnan; <his was  not a bad dream! What a speculation! Peste! Well! what is said is said.>

<How much are you to give him?>

<Oh!> said D'Artagnan, <he is not a bad fellow; I shall  arrange matters with him. I have had a great deal of trouble, you see, and  expenses; all that must be taken into account.>

<My dear friend, I can depend on you, and have no fear for the worthy  Planchet; his interests are better in your hands than in his own. But now that  you have nothing more to do here, we shall depart, if you please. You can go  and thank his majesty, ask if he has any commands, and in six days we may be  able to get sight of the towers of Notre Dame.>

<My friend, I am most anxious to be off, and will go at once and pay my  respects to the king.>

<I,> said Athos, <am going to call upon some friends in the  city, and shall then be at your service.>

<Will you lend me Grimaud?>

<With all my heart. What do you want to do with him?>

<Something very simple, and which will not fatigue him; I shall only beg  him to take charge of my pistols, which lie there on the table near that  coffer.>

<Very well!> replied Athos, imperturbably.

<And he will not stir, will he?>

<Not more than the pistols themselves.>

<Then I shall go and take leave of his majesty. Au revoir!>

D'Artagnan arrived at St James's, where Charles II, who was busy  writing, kept him in the ante-chamber a full hour. Whilst walking about in the  gallery, from the door to the window, from the window to the door, he thought  he saw a cloak like Athos's cross the vestibule; but at the moment he was  going to ascertain if it were he, the usher summoned him to his majesty's  presence. Charles II rubbed his hands while receiving the thanks of our  friend.

<Chevalier,> said he, <you are wrong to express gratitude to  me; I have not paid you a quarter of the value of the history of the box into  which you put the brave general—the excellent Duke of Albemarle, I  mean.> And the king laughed heartily.

D'Artagnan did not think it proper to interrupt his majesty, and he bowed  with much modesty.

<A propos,> continued Charles, <do you think my dear Monk has  really pardoned you?>

<Pardoned me! yes, I hope so, sire!>

<Eh!—but it was a cruel trick! Odds fish! to pack up the first  personage of the English revolution like a herring. In your place I would not  trust him, chevalier.>

<But, sire\longdash>

<Yes, I know very well Monk calls you his friend, but he has too  penetrating an eye not to have a memory, and too lofty a brow not to be very  proud, you know, grande supercilium.>

<I shall certainly learn Latin,> said D'Artagnan to himself.

<But stop,> cried the merry monarch, <I must manage your  reconciliation; I know how to set about it; so\longdash>

D'Artagnan bit his moustache. <Will your majesty permit me to tell  you the truth?>

<Speak, chevalier, speak.>

<Well, sire, you alarm me greatly. If your majesty undertakes the affair,  as you seem inclined to do, I am a lost man; the duke will have me  assassinated.>

The king burst into a fresh roar of laughter, which changed  D'Artagnan's alarm into downright terror.

<Sire, I beg you to allow me to settle this matter myself, and if your  majesty has no further need of my services\longdash>

<No, chevalier. What, do you want to leave us?> replied Charles,  with a hilarity that grew more and more alarming.

<If your majesty has no more commands for me.>

Charles became more serious.

<One single thing. See my sister, the Lady Henrietta. Do you know  her?>

<No, sire, but—an old soldier like me is not an agreeable spectacle  for a young and gay princess.>

<Ah! but my sister must know you; she must in case of need have you to  depend upon.>

<Sire, every one that is dear to your majesty will be sacred to  me.>

<Very well!—Parry! Come here, Parry!>

The side door opened and Parry entered, his face beaming with pleasure as soon  as he saw D'Artagnan.

<What is Rochester doing?> said the king.

<He is on the canal with the ladies,> replied Parry.

<And Buckingham?>

<He is there also.>

<That is well. You will conduct the chevalier to Villiers; that is the  Duke of Buckingham, chevalier; and beg the duke to introduce M.  d'Artagnan to the Princess Henrietta.>

Parry bowed and smiled to D'Artagnan.

<Chevalier,> continued the king, <this is your parting  audience; you can afterwards set out as soon as you please.>

<Sire, I thank you.>

<But be sure you make your peace with Monk!>

<Oh, sire\longdash>

<You know there is one of my vessels at your disposal?>

<Sire, you overpower me; I cannot think of putting your majesty's  officers to inconvenience on my account.>

The king slapped D'Artagnan upon the shoulder.

<Nobody will be inconvenienced on your account, chevalier, but for that  of an ambassador I am about sending to France, and to whom you will willingly  serve as a companion, I fancy, for you know him.>

D'Artagnan appeared astonished.

<He is a certain Comte de la Fère,—whom you call Athos,>  added the king; terminating the conversation, as he had begun it, by a joyous  burst of laughter. <Adieu, chevalier, adieu. Love me as I love  you.> And thereupon, making a sign to Parry to ask if there were any one  waiting for him in the adjoining closet, the king disappeared into that closet,  leaving the chevalier perfectly astonished by this singular audience. The old  man took his arm in a friendly way, and led him towards the garden.  